\chapter{Overview}

This blueprint documents the constructive $P(\Phi)_2$ Euclidean quantum field
theory in the \texttt{pphi2} repository together with the upstream results it
uses from our own libraries (\texttt{gaussian-field}).

\emph{Edges are extracted, not guessed.} Every \texttt{\textbackslash uses} edge
below is computed from the \emph{actual} Lean dependency graph
(\texttt{scripts/BlueprintDeps.lean} walks each declaration's transitive
used-constant closure, restricted to our namespaces, then transitively reduces
it). Green nodes have a formalized proof; the node flagged \textbf{axiom} is a
Lean \texttt{axiom} (assumed, not proved). Nodes marked \textbf{(upstream)} live
in \texttt{gaussian-field}.

\emph{A structural finding.} The dependency closure of the headline theorem
\texttt{pphi2\_exists} is in fact \emph{small}: it bottoms out in the OS-axiom
inheritance lemmas and the clustering \textbf{axiom}. Several substantial
\emph{proved} results --- tightness, the transfer-matrix mass gap,
reflection positivity from the lattice, lattice translation invariance, and the
Ward/anomaly analysis --- are developed in the repository but are \emph{not}
(yet) part of \texttt{pphi2\_exists}'s proof closure. They appear here in
Chapter~\ref{chap:supporting} as standalone components, which is the honest
picture of the current formalization.

\chapter{Main theorem and the OS axioms}

\begin{definition}[Full OS axioms]\label{def:full-os}
  \lean{Pphi2.SatisfiesFullOS}\leanok
  Bundles OS0 (analyticity), OS1 (regularity), OS2 (Euclidean invariance),
  OS3 (reflection positivity), OS4 (clustering + ergodicity).
\end{definition}

\begin{theorem}[Continuum exponential moments]\label{thm:expmom}
  \lean{Pphi2.continuum_exponential_moments}\leanok
  The continuum-limit measure has finite exponential moments.
\end{theorem}

\begin{theorem}[OS0 in the limit]\label{thm:os0}
  \lean{Pphi2.os0_for_continuum_limit}\leanok\uses{thm:expmom}
\end{theorem}
\begin{theorem}[OS1 in the limit]\label{thm:os1}
  \lean{Pphi2.os1_for_continuum_limit}\leanok\uses{thm:expmom}
\end{theorem}

\begin{theorem}[Real generating functional invariance]\label{thm:gfinv}
  \lean{Pphi2.complex_gf_invariant_of_real_gf_invariant}\leanok
  Real Euclidean invariance of the generating functional upgrades to the
  complex statement (given exponential moments).
\end{theorem}
\begin{theorem}[Translation invariance (continuum)]\label{thm:transinv}
  \lean{Pphi2.translation_invariance_continuum}\leanok
\end{theorem}
\begin{theorem}[Rotation invariance (continuum)]\label{thm:rotinv}
  \lean{Pphi2.rotation_invariance_continuum}\leanok\uses{def:gff-measure}
\end{theorem}
\begin{theorem}[OS2 (Euclidean invariance) in the limit]\label{thm:os2}
  \lean{Pphi2.os2_for_continuum_limit}\leanok
  \uses{thm:gfinv, thm:expmom, thm:rotinv, thm:transinv}
\end{theorem}

\begin{theorem}[OS3 (reflection positivity) in the limit]\label{thm:os3}
  \lean{Pphi2.os3_for_continuum_limit}\leanok
\end{theorem}

\begin{theorem}[Continuum exponential clustering \textbf{(axiom --- assumed)}]\label{ax:clustering}
  \lean{Pphi2.continuum_exponential_clustering}\leanok
  Exponential decay of the connected two-point function. \emph{Currently a Lean
  \texttt{axiom}.}
\end{theorem}
\begin{theorem}[OS4 clustering in the limit]\label{thm:os4}
  \lean{Pphi2.os4_for_continuum_limit}\leanok\uses{ax:clustering}
\end{theorem}
\begin{theorem}[Clustering $\Rightarrow$ ergodicity]\label{thm:os4-erg}
  \lean{Pphi2.os4_clustering_implies_ergodicity}\leanok
\end{theorem}

\begin{theorem}[The limit satisfies OS0--OS4]\label{thm:limit-os}
  \lean{Pphi2.continuumLimit_satisfies_fullOS}\leanok
  \uses{def:full-os, thm:os0, thm:os1, thm:os2, thm:os3, thm:os4, thm:os4-erg}
\end{theorem}

\begin{theorem}[Continuum limit exists]\label{thm:limit-exists}
  \lean{Pphi2.pphi2_limit_exists}\leanok
  The lattice $P(\Phi)_2$ measures have a continuum-limit probability measure
  (tightness + Prokhorov; the latter axiomatized in Mathlib terms).
\end{theorem}

\begin{theorem}[Existence of a $P(\Phi)_2$ OS measure]\label{thm:main}
  \lean{Pphi2.pphi2_exists}\leanok\uses{thm:limit-os, thm:limit-exists}
  For every interaction polynomial $P$ and mass $m>0$ there is a probability
  measure on cylinder configurations satisfying the full OS axioms.
\end{theorem}

\chapter{Upstream Gaussian measure (\texttt{gaussian-field})}

\begin{definition}[Gaussian field measure (upstream)]\label{def:gff-measure}
  \lean{GaussianField.measure}\leanok
  The centered Gaussian measure for a continuous linear map $T : E \to H$, from
  its characteristic functional.
\end{definition}
\begin{definition}[Covariance form (upstream)]\label{def:gff-covariance}
  \lean{GaussianField.covariance}\leanok
  $C_T(f,g) = \langle Tf, Tg\rangle_H$.
\end{definition}

\chapter{Supporting proved results (not yet in the headline closure)}\label{chap:supporting}

These are formalized (green) but are \emph{not} currently in the dependency
closure of Theorem~\ref{thm:main}; wiring them in is future work.

\section{Tightness}
\begin{theorem}[Configuration tightness (upstream)]\label{thm:tight-upstream}
  \lean{GaussianField.configuration_tight_of_uniform_second_moments}\leanok
\end{theorem}
\begin{theorem}[The lattice family is tight]\label{thm:tight}
  \lean{Pphi2.continuumMeasures_tight}\leanok
  \uses{thm:tight-upstream, def:gff-covariance, def:gff-measure}
\end{theorem}

\section{Transfer-matrix mass gap}
\begin{theorem}[Energy gap]\label{thm:energygap}
  \lean{Pphi2.energyLevel_gap}\leanok
\end{theorem}
\begin{theorem}[Lattice mass gap positive]\label{thm:massgap}
  \lean{Pphi2.massGap_pos}\leanok\uses{thm:energygap}
\end{theorem}
\begin{theorem}[OS4 from the gap (lattice)]\label{thm:os4gap}
  \lean{Pphi2.os4_lattice_from_gap}\leanok\uses{thm:massgap}
\end{theorem}

\section{Reflection positivity from the lattice}
\begin{theorem}[Time reflection is an involution]\label{thm:timerefl}
  \lean{Pphi2.timeReflection2D_involution}\leanok
\end{theorem}
\begin{theorem}[Action reflection decomposition]\label{thm:actiondecomp}
  \lean{Pphi2.action_decomposition}\leanok\uses{thm:timerefl}
\end{theorem}
\begin{theorem}[RP closed under weak limits]\label{thm:rpweak}
  \lean{Pphi2.rp_closed_under_weak_limit}\leanok
\end{theorem}
\begin{theorem}[Continuum RP from lattice RP]\label{thm:rpcont}
  \lean{Pphi2.continuum_rp_from_lattice}\leanok\uses{thm:rpweak}
\end{theorem}

\section{Lattice translation invariance}
\begin{theorem}[Mass operator commutes with translation]\label{thm:massop}
  \lean{Pphi2.massOperator_translation_commute}\leanok
\end{theorem}
\begin{theorem}[Gaussian density translation invariant]\label{thm:gdens}
  \lean{Pphi2.gaussianDensity_translation_invariant}\leanok\uses{thm:massop}
\end{theorem}
\begin{theorem}[Lattice action translation invariant]\label{thm:latact}
  \lean{Pphi2.latticeAction_translation_invariant}\leanok
\end{theorem}
\begin{theorem}[Lattice measure translation invariant]\label{thm:latmeas}
  \lean{Pphi2.latticeMeasure_translation_invariant}\leanok
  \uses{def:gff-covariance, def:gff-measure, thm:gdens}
\end{theorem}

\section{Ward identity and anomaly}
\begin{theorem}[Lattice Ward identity]\label{thm:ward}
  \lean{Pphi2.ward_identity_lattice}\leanok\uses{def:gff-measure}
\end{theorem}
\begin{theorem}[Anomaly bound from super-renormalizability]\label{thm:anombound}
  \lean{Pphi2.anomaly_bound_from_superrenormalizability}\leanok\uses{def:gff-measure}
\end{theorem}
\begin{theorem}[Rotational CF defect bound]\label{thm:cfdefect}
  \lean{Pphi2.rotationCFDefect_le_pointwiseDefectIntegral}\leanok\uses{def:gff-measure}
\end{theorem}
\begin{theorem}[Anomaly vanishes]\label{thm:anomvanish}
  \lean{Pphi2.anomaly_vanishes}\leanok\uses{thm:anombound}
\end{theorem}
